%% kapa: exam=1 round=26 year=2015 period=1 form=A subject=civil_law subject_ko=민법 seq=32 number=32
%% source: Q-Net 감정평가사 기출문제 (cst00309) artlSeq=5138001 file=extracted/1008/5138001/제26회_감정평가사 1차_기출/제26회_감정평가사 1차_기출/1교시_감평_A형.pdf
%% conversion: kordoc 4.24.51 -> tools/gen1.py (auto); review=proofread-sonnet; flags=none
\begin{question}{32}
물권변동에 관한 설명으로 옳지 않은 것은? (다툼이 있으면 판례에 따름)
\begin{choices}
  \item 수공업자가 타인의 동산에 가공을 한 경우 가공으로 인한 가액의 증가가 원재료의 가액보다 현저히 다액인 경우, 가공물은 그 수공업자의 소유로 한다.
  \item 부동산등기법상 무효인 이중등기를 근거로 해서도 등기부취득시효가 완성될 수 있다.
  \item 인도는 법률행위로 인한 부동산 소유권취득을 위한 요건이 아니다.
  \item 타인의 임야에 권원 없이 심은 입목의 소유권은 임야소유자에게 속한다.
  \item 경매 대상 토지에 채무자 아닌 자의 소유인 동산이 설치되어 있는 경우, 경매의 매수인은 그 동산을 선의취득할 수 있다.
\end{choices}
\end{question}
\answer{2}
